% ================================================================
% 5. PROCESSI ORGANIZZATIVI
% ================================================================
\section{Processi Organizzativi}
 
\subsection{Processo Di Gestione}
\label{sec:gestione}
 
\subsubsection{Descrizione}
 
Il processo di gestione regola la pianificazione delle attività, l'assegnazione e
la rotazione dei ruoli, la conduzione delle riunioni e il tracciamento
dell'avanzamento del lavoro. Esso fornisce la struttura organizzativa entro cui
tutti gli altri processi operano.
 
\subsubsection{Ruoli Del Gruppo}
\label{sec:ruoli}
 
Il gruppo prevede i seguenti ruoli, assegnati a rotazione tra i membri. Il
Regolamento del Progetto Didattico impone che ogni componente:
 
\begin{itemize}
  \item ruoti su tutti i ruoli secondo regole documentate;
  \item assuma ogni ruolo per un tempo totale congruo;
  \item assuma \textbf{non più di un ruolo alla volta}.
\end{itemize}
 
Le regole di rotazione devono assicurare assenza di conflitti di interesse,
garantendo verifica indipendente per ogni attività.
 
\bigskip
\begin{tabularx}{\textwidth}{l c X}
  \toprule
  \textbf{Ruolo} & \textbf{Costo orario} & \textbf{Responsabilità} \\
  \midrule
  Responsabile\textsubscript{G}
    & 30\,€/h
    & Coordina l'elaborazione di piani e scadenze; approva il rilascio di
      prodotti parziali o finali (SW, documenti); coordina le attività del
      gruppo; rappresenta \textit{Synergo Unit} verso committente e
      proponente. \\[6pt]
  Amministratore\textsubscript{G}
    & 20\,€/h
    & Assicura l'efficienza di procedure, strumenti e tecnologie a supporto
      del \textit{way of working}; redige e mantiene aggiornate le Norme di
      Progetto. \\[6pt]
  Analista\textsubscript{G}
    & 25\,€/h
    & Svolge le attività di analisi dei requisiti; attivato dopo
      l'aggiudicazione del capitolato. \\[6pt]
  Progettista
    & 25\,€/h
    & Svolge le attività di progettazione (\textit{design}); attivato dopo
      l'aggiudicazione del capitolato. \\[6pt]
  Programmatore\textsubscript{G}
    & 15\,€/h
    & Svolge le attività di codifica; attivato durante la fase di sviluppo. \\[6pt]
  Verificatore\textsubscript{G}
    & 15\,€/h
    & Svolge le attività di verifica di documenti e codice prodotti dagli
      altri membri; redige i relativi esiti di verifica. \\
  \bottomrule
\end{tabularx}
 
\bigskip
Le assegnazioni di ruolo sono tracciate nel documento
\textit{Dichiarazione Degli Impegni} e nei task\textsubscript{G} su
GitHub Project\textsubscript{G}.
 
\subsubsection{Pianificazione E Tracciamento Delle Attività}
 
Le attività sono tracciate tramite \textbf{GitHub Project\textsubscript{G}}.
Ogni task\textsubscript{G} è identificato dal codice \texttt{T-y.z} e riporta
obbligatoriamente i seguenti attributi:
 
\begin{itemize}
  \item assegnatario e ruolo;
  \item decisione di riferimento (codice \texttt{VI-y.x} o \texttt{VE-y.x});
  \item data di scadenza;
  \item stima del tempo previsto (in minuti);
  \item tempo effettivo consuntivato (in minuti);
  \item stato di avanzamento (\textit{To Do}, \textit{In Progress},
        \textit{Done}).
\end{itemize}
 
La creazione di un task avviene tramite \textbf{Google Form}, che genera
automaticamente una issue\textsubscript{G} su GitHub e registra i dati in un
foglio di calcolo condiviso, utilizzato per il monitoraggio aggregato delle
misurazioni.
 
\subsubsection{Riunioni Interne}
 
\paragraph{Cadenza E Piattaforma}
Il gruppo si riunisce con cadenza fissa, salvo diversa comunicazione del
Responsabile:
 
\begin{itemize}
  \item \textbf{Martedì} alle \textbf{14:30};
  \item \textbf{Venerdì} alle \textbf{14:30}.
\end{itemize}
 
Le riunioni si svolgono in modalità remota sulla piattaforma Discord\textsubscript{G}.
Il Responsabile può convocare riunioni straordinarie tramite il gruppo
WhatsApp\textsubscript{G}.
 
\paragraph{Struttura Della Riunione}
Ogni riunione segue la struttura seguente:
 
\begin{enumerate}
  \item \textbf{Predisposizione dell'ordine del giorno}: il Responsabile pubblica
        i punti da discutere su un documento Google condiviso, nel quale vengono
        anche registrate le decisioni prese; tale documento costituisce la base
        per la successiva redazione del verbale.
  \item \textbf{Discussione dei punti all'ordine del giorno}: ogni membro può
        contribuire; le decisioni vengono adottate collegialmente.
  \item \textbf{Assegnazione dei task}: al termine della discussione vengono
        definiti e assegnati i task conseguenti, con scadenza e ruolo.
  \item \textbf{Redazione del verbale}: il Redattore incaricato compila il verbale
        interno al termine della riunione; il verbale è verificato e pubblicato
        entro la riunione successiva.
\end{enumerate}
 
\paragraph{Modalità Decisionale}
Il gruppo adotta il \textbf{brainstorming} come tecnica principale per la generazione
e il confronto delle idee. Il processo si articola in tre fasi:
 
\begin{enumerate}
  \item raccolta libera delle proposte da parte di ogni membro, senza valutazione
        immediata;
  \item confronto strutturato delle idee emerse, con analisi di vantaggi e
        svantaggi;
  \item scelta collegiale della soluzione ritenuta più adeguata, con registrazione
        della decisione nel verbale.
\end{enumerate}
 
Le decisioni adottate sono registrate con il codice \texttt{VI-y.x}. Le decisioni
non possono essere modificate unilateralmente: eventuali revisioni richiedono
delibera in una riunione successiva.
 
\subsubsection{Gestione Delle Disponibilità}
 
Le disponibilità dei membri e la pianificazione delle riunioni ricorrenti sono
gestite tramite \textbf{Google Calendar\textsubscript{G}}. I calendari condivisi
sono aggiornati dal Responsabile o dall'Amministratore. Le comunicazioni urgenti
e le notifiche rapide avvengono tramite il gruppo \textbf{WhatsApp\textsubscript{G}}.
 
% ---------------------------------------------------------------
\subsection{Processo Di Infrastruttura}
\label{sec:infrastruttura}
 
\subsubsection{Descrizione}
 
Il processo di infrastruttura definisce, mantiene e documenta l'insieme degli
strumenti adottati dal gruppo a supporto di tutti i processi attivi. L'adozione
di un nuovo strumento deve essere deliberata collegialmente e registrata nel verbale
della riunione in cui viene approvata; le presenti norme devono essere aggiornate
di conseguenza.
 
\subsubsection{Strumenti Di Gestione}
 
\begin{itemize}
  \item \textbf{GitHub Project\textsubscript{G}}: sistema principale per la
        pianificazione e il tracciamento dei task\textsubscript{G}, con gestione
        di assegnatari, scadenze, stime e stati di avanzamento.
  \item \textbf{Google Form}: interfaccia guidata per la creazione dei task;
        genera automaticamente la issue\textsubscript{G} corrispondente su GitHub
        e registra i dati nel foglio di calcolo condiviso.
  \item \textbf{Foglio Di Calcolo}: aggregazione e monitoraggio delle misurazioni
        relative ai task (tempo previsto, tempo effettivo, stato di avanzamento).
  \item \textbf{Google Calendar\textsubscript{G}}: gestione delle disponibilità
        dei membri e pianificazione delle riunioni ricorrenti e straordinarie.
\end{itemize}
 
\subsubsection{Strumenti Di Documentazione}
 
\begin{itemize}
  \item \textbf{\LaTeX{} con Visual Studio Code}: ambiente di redazione documentale
        principale; la compilazione automatica in PDF è gestita tramite l'estensione
        \textit{LaTeX Workshop}.
  \item \textbf{Script Python} (\texttt{python\_glossary\_automation/}):
        automatizza l'applicazione del pedice G\textsubscript{G} ai termini del
        Glossario, riducendo il rischio di omissioni.
  \item \textbf{GitHub}: sistema di controllo versione centralizzato per tutti i
        documenti del progetto; gestisce branch\textsubscript{G}, pull
        request\textsubscript{G} e baseline\textsubscript{G}.
  \item \textbf{GitHub Pages\textsubscript{G}}: piattaforma di pubblicazione del
        sito documentale del gruppo, sviluppato in HTML5 e CSS3.
\end{itemize}
 
\subsubsection{Strumenti Di Comunicazione}
 
\begin{itemize}
  \item \textbf{Email Istituzionale Del Gruppo}: canale ufficiale per le
        comunicazioni formali con committente e proponente. Configurata con
        inoltro automatico verso i membri delegati; le risposte mantengono
        l'indirizzo istituzionale come mittente.
  \item \textbf{Discord\textsubscript{G}}: piattaforma principale per le riunioni
        interne, le discussioni operative e il coordinamento quotidiano tra i
        membri.
  \item \textbf{WhatsApp\textsubscript{G}}: canale secondario per comunicazioni
        rapide, notifiche urgenti e convocazioni straordinarie.
\end{itemize}
 